%% kapa: exam=1 round=33 year=2022 period=2 form=A subject=accounting subject_ko=회계학 seq=6 number=46
%% source: Q-Net 감정평가사 기출문제 (cst00309) artlSeq=5209997 file=raw/1008/5209997/2022년도 제33회 감정평가사 1차시험 2교시(감정평가 관계법규, 회계학).pdf
%% conversion: kordoc 4.24.51 -> tools/gen1.py (auto); review=proofread-sonnet; flags=none
\begin{question}{46}
(주)감평의 20×1년 기말 재고자산 자료가 다음과 같다.
\par{\small\centering\begin{tabular}{|l|r|r|r|}\hline
종 목 & 실사수량 & 단위당 취득원가 & 단위당 예상판매가격 \\ \hline
상품 A & 100개 & ￦300 & ￦350 \\
상품 B & 100개 & 200 & 250 \\
상품 C & 200개 & 100 & 120 \\ \hline
\end{tabular}\par}
\begin{framedtext}
○ 단위당 예상판매비용: ￦30(모든 상품에서 발생) \par
\end{framedtext}
상품 B의 70\%는 확정판매계약(취소불능계약)을 이행하기 위하여 보유하고 있으며, 상품 B의 단위당 확정판매계약가격은 ￦220이다. 재고자산 평가와 관련하여 20×1년 인식할 당기손익은? (단, 재고자산의 감모는 발생하지 않았으며, 기초 재고자산평가충당금은 없다.)
\begin{choices}
  \item 손실 ￦2,700
  \item 손실 ￦700
  \item ￦0
  \item 이익 ￦2,200
  \item 이익 ￦3,200
\end{choices}
\end{question}
\answer{1}
